\documentclass[11pt,a4paper]{article}
\usepackage[utf8]{inputenc}
\usepackage{amsmath,amssymb,amsthm}
\usepackage{algorithm}
\usepackage{algpseudocode}
\usepackage{booktabs}
\usepackage{hyperref}
\usepackage{geometry}
\geometry{margin=1in}

\title{Sequence-to-Sequence (Seq2Seq) Encoder-Decoder Networks: Information Bottleneck Mechanics and Autoregressive Decoding}
\author{Team Scaibu \\ \small Email: \texttt{jaydeep.9664920749@gmail.com} \\ \small Architecture and Core Engine Team}
\date{\today}

\newtheorem{theorem}{Theorem}
\newtheorem{lemma}[theorem]{Lemma}
\newtheorem{proposition}[theorem]{Proposition}
\theoremstyle{definition}
\newtheorem{definition}{Definition}
\newtheorem{remark}{Remark}

\begin{document}

\maketitle

\begin{abstract}
Traditional neural networks assume fixed-dimensional mappings between input and output representations, failing when sequence lengths vary dynamically ($T_x \ne T_y$) as in machine translation, text summarization, and query synthesis. The Sequence-to-Sequence (Seq2Seq) architecture resolves this asymmetry by decoupling the encoding of an arbitrary-length source sequence from the autoregressive conditional decoding of the target sequence. This document formalizes the information-theoretic vector bottleneck between encoder and decoder, proves maximum likelihood factorization properties, details algorithmic pseudocode, and provides a verified numerical walkthrough.
\end{abstract}

\section{Notation and Mathematical Preliminaries}

Let $\mathbf{x} = (x_1, \dots, x_{T_x})$ denote a source sequence of token indices from vocabulary $\mathcal{V}_x$, and let $\mathbf{y} = (y_1, \dots, y_{T_y})$ denote a target sequence from vocabulary $\mathcal{V}_y$.

\begin{table}[h]
\centering
\caption{Mathematical Notation for Seq2Seq Architecture}
\begin{tabular}{lll}
\toprule
\textbf{Symbol} & \textbf{Domain} & \textbf{Semantic Description} \\
\midrule
$\mathbf{h}_t^{(\text{enc})}$ & $\mathbb{R}^{d_h}$ & Encoder hidden state vector at input step $t \in [1, T_x]$ \\
$\mathbf{v}$ & $\mathbb{R}^{d_h}$ & Fixed-dimensional source context vector: $\mathbf{v} = \mathbf{h}_{T_x}^{(\text{enc})}$ \\
$\mathbf{s}_u^{(\text{dec})}$ & $\mathbb{R}^{d_h}$ & Decoder hidden state vector at output step $u \in [1, T_y]$ ($\mathbf{s}_0 = \mathbf{v}$) \\
$E_{\text{src}}, E_{\text{tgt}}$ & $\mathbb{R}^{|\mathcal{V}| \times d_e}$ & Source and target token embedding matrices \\
$W_{\text{vocab}}, \mathbf{b}_{\text{vocab}}$ & Matrix, Vector & Target vocabulary projection emission parameters \\
$\langle \text{BOS} \rangle, \langle \text{EOS} \rangle$ & Tokens & Special boundary tokens indicating beginning and end of sequence \\
\bottomrule
\end{tabular}
\end{table}

\section{Problem Definition and Core Concepts}

\subsection{3-Level Explanation}
\begin{enumerate}
    \item \textbf{Intuitive (High School level):} Imagine a translator listening to a complete sentence in French, forming a mental concept of its meaning (Encoder), and then speaking out the English translation one word at a time until the thought is complete (Decoder).
    \item \textbf{Engineering Level:} An Encoder RNN consumes source embeddings $E(x_t)$, accumulating state into a final context vector $\mathbf{v} = \mathbf{h}_{T_x}$. The Decoder RNN initializes its initial state with $\mathbf{s}_0 = \mathbf{v}$, conditioning each next token prediction on past emitted tokens $\mathbf{y}_{<u}$ until emitting the stop token $\langle \text{EOS} \rangle$.
    \item \textbf{Mathematical Level:} Seq2Seq computes the exact joint probability distribution of the target sequence conditioned on the source via the chain rule of probability:
    \begin{equation}
    P(\mathbf{y} \mid \mathbf{x}) = \prod_{u=1}^{T_y} P(y_u \mid y_1, \dots, y_{u-1}, \mathbf{x}) = \prod_{u=1}^{T_y} \text{Softmax}\left( W_{\text{vocab}} \mathbf{s}_u^{(\text{dec})} + \mathbf{b}_{\text{vocab}} \right)_{y_u}
    \end{equation}
\end{enumerate}

\subsection{5-Step Algorithmic Framework}
\begin{enumerate}
    \item \textbf{Source Sequence Embedding:} Convert token IDs $x_1, \dots, x_{T_x}$ into continuous embeddings.
    \item \textbf{Encoder State Accumulation:} Process embeddings sequentially via encoder RNN to extract context vector $\mathbf{v} = \mathbf{h}_{T_x}^{(\text{enc})}$.
    \item \textbf{Decoder Initialization:} Set initial decoder state $\mathbf{s}_0 = \mathbf{v}$ and feed start token $y_0 = \langle \text{BOS} \rangle$.
    \item \textbf{Autoregressive Step Emission:} For step $u = 1 \dots T_{\max}$, compute $\mathbf{s}_u = \text{RNN}(\mathbf{s}_{u-1}, E(y_{u-1}))$ and emit token $y_u = \arg\max_k P(k \mid \mathbf{s}_u)$.
    \item \textbf{EOS Termination:} Stop generation if $y_u = \langle \text{EOS} \rangle$ or when reaching length limit.
\end{enumerate}

\section{Algorithm Specification}

\begin{algorithm}[H]
\caption{Seq2Seq Greedy Autoregressive Generation}
\label{alg:seq2seq}
\begin{algorithmic}[1]
\Require Source token sequence $\mathbf{x} = (x_1, \dots, x_{T_x})$, embedding tables $E_{\text{src}}, E_{\text{tgt}}$, encoder parameters $\theta_{\text{enc}}$, decoder parameters $\theta_{\text{dec}}$, max length $T_{\max}$.
\Ensure Decoded sequence $\mathbf{y} = (y_1, \dots, y_{T_y})$.
\State $\mathbf{h}_0 \gets \mathbf{0}$
\For{$t = 1$ \textbf{to} $T_x$} \Comment{Encoder Pass}
    \State $\mathbf{h}_t \gets \tanh(W_h \mathbf{h}_{t-1} + W_x E_{\text{src}}[x_t] + \mathbf{b}_h)$
\EndFor
\State $\mathbf{s}_0 \gets \mathbf{h}_{T_x}, \quad y_0 \gets \langle \text{BOS} \rangle, \quad \mathbf{y} \gets [\,]$
\For{$u = 1$ \textbf{to} $T_{\max}$} \Comment{Autoregressive Decoder Rollout}
    \State $\mathbf{s}_u \gets \tanh(W_s \mathbf{s}_{u-1} + W_y E_{\text{tgt}}[y_{u-1}] + \mathbf{b}_s)$
    \State $\mathbf{z}_u \gets W_{\text{vocab}} \mathbf{s}_u + \mathbf{b}_{\text{vocab}}$
    \State $y_u \gets \arg\max_{k} \mathbf{z}_u[k]$
    \If{$y_u = \langle \text{EOS} \rangle$}
        \State \textbf{break}
    \EndIf
    \State Append $y_u$ to $\mathbf{y}$
\EndFor
\State \Return $\mathbf{y}$
\end{algorithmic}
\end{algorithm}

\section{Theoretical Guarantees and Information Bottleneck Bound}

\begin{theorem}[Vector Bottleneck Capacity Bound]
\label{thm:bottleneck}
Let the context vector be bounded in dimension $\mathbf{v} \in \mathbb{R}^{d_h}$. The mutual information between the source sequence $\mathbf{x}$ and the generated translation $\mathbf{y}$ is strictly upper-bounded by the capacity of the continuous hidden vector:
\begin{equation}
I(\mathbf{x}; \mathbf{y}) \le I(\mathbf{x}; \mathbf{v}) \le d_h \cdot \log_2\left(1 + \frac{\text{Var}(\mathbf{v})}{\sigma_{\text{noise}}^2}\right)
\end{equation}
proving that fixed-vector Seq2Seq suffers from unavoidable information degradation on long inputs ($T_x \gg d_h$).
\end{theorem}

\begin{proof}
Forming the Markov chain $\mathbf{x} \to \mathbf{v} \to \mathbf{s}_u \to \mathbf{y}$, the Data Processing Inequality dictates $I(\mathbf{x}; \mathbf{y}) \le I(\mathbf{x}; \mathbf{v})$. For a $d_h$-dimensional continuous Gaussian channel with signal variance $\text{Var}(\mathbf{v})$ and numerical quantization noise $\sigma_{\text{noise}}^2$, Shannon channel capacity bounds maximum mutual information by $d_h \log_2(1 + \text{SNR})$.
\end{proof}

\section{Complexity Analysis}

\begin{itemize}
    \item \textbf{Time Complexity:} $\mathcal{O}(T_x \cdot d_h^2 + T_y \cdot (d_h^2 + d_h |\mathcal{V}|))$ operations (strictly sequential during inference).
    \item \textbf{Space Complexity:} $\mathcal{O}(d_h)$ working memory for the hidden state vector.
\end{itemize}

\section{Exactness and Error Analysis}

Greedy autoregressive decoding searches only a single trajectory, potentially missing the globally optimal sequence $\mathbf{y}^* = \arg\max_\mathbf{y} P(\mathbf{y} \mid \mathbf{x})$. Beam search with beam width $B$ expands the search space to approximate optimal decoding.

\section{Worked Numerical Example}

Let $d_h=1, d_e=1, |\mathcal{V}|=3$ ($0=\text{PAD}, 1=\text{BOS}, 2=\text{EOS}$).
Let encoder produce context vector $v = 0.5$.
Decoder parameters: $W_s = [1.0], W_y = [1.0], b_s = [0.0]$, $W_{\text{vocab}} = [[-1.0], [0.0], [2.0]]$.
At step $u=1$ with $y_0 = 1$ (emb $= 0.0$):
\begin{itemize}
    \item $s_1 = \tanh(1.0(0.5) + 1.0(0.0)) = \tanh(0.5) = 0.4621$.
    \item Logits: $z_0 = -0.4621, z_1 = 0.0, z_2 = 2.0(0.4621) = 0.9242$.
    \item $\arg\max_k z = 2 = \langle \text{EOS} \rangle$. Generation halts.
\end{itemize}

\section{Numerical Considerations and Edge Cases}

\begin{itemize}
    \item \textbf{Reversing Source Sequence:} In pure RNN Seq2Seq, reversing the source sequence ($x_{T_x}, \dots, x_1$) dramatically improves performance by shortening the minimal time lag between early source tokens and early target tokens.
\end{itemize}

\section{References}
\begin{enumerate}
    \item Sutskever, I., Vinyals, O., & Le, Q. V. (2014). Sequence to sequence learning with neural networks. \emph{NeurIPS}, 3104--3112.
    \item Cho, K., Van Merri{\"e}nboer, B., Gulcehre, C., Bahdanau, D., Bougares, F., Schwenk, H., & Bengio, Y. (2014). Learning phrase representations using RNN encoder-decoder for statistical machine translation. \emph{EMNLP}, 1724--1734.
\end{enumerate}

\end{document}
